\subsection{不可约元素}

定义 6.——序群 G 的元素 x 称为不可约的，如果它是 G 的严格正元素集合中的极小元素。

设 x 是序群 G 的一个不可约元；若 y 是 G 的正元素，则元素 \(\inf(x, y)\) 若存在，只能等于 x 或 0。因此在格序群 G 中，每个正元素 y 或者大于不可约元 x，或者与 x 互素；特别地，两个不同的不可约元是互素的。

\begin{enumerate}
\def\labelenumi{(\Roman{enumi})}
\setcounter{enumi}{503}
\tightlist
\item
  元素 \(p\)（属于 \(\mathbf{A}\)）称为不可约的，如果理想 \((p)\) 是序群 \(\mathcal{P}^*\) 的一个不可约元；这表示 \(p\) 既非零亦非可逆元，且 \(\mathbf{A}\) 中任意整除 \(p\) 的元素与 \(p\) 或与 1 相伴。如果 \(\mathcal{P}^*\) 是格序的，那么每个 \(a \in \mathbf{A}\) 要么与 \(p\) 互素，要么是 \(p\) 的倍元。
\end{enumerate}

例子（DIV）。—— 1）整数 \(p > 0\) 在 \(\mathbf{Z}\) 中不可约，当且仅当它是素数（I，第 50 页）。

\begin{enumerate}
\def\labelenumi{\arabic{enumi})}
\setcounter{enumi}{1}
\tightlist
\item
  域 K 上关于单个未定元的多项式是不可约的，当且仅当它在通常意义下不可约（IV，第 13 页）。
\end{enumerate}

命题 14.——要使 \(x > 0\) （序群 \(G\) 的元素）成为不可约元，只需它满足以下性质：

\begin{enumerate}
\def\labelenumi{(\Alph{enumi})}
\setcounter{enumi}{15}
\tightlist
\item
  关系 \(x \leqslant y + z\) , \(y \geqslant 0\) , \(z \geqslant 0\) 蕴含 \(x \leqslant y\) 或 \(x \leqslant z\) 。
\end{enumerate}

当 G 是格序群时，这个条件是必要的。

如果 G 是格序的且 x 不可约，我们刚刚看到，y 要么大于 x，要么与 x 互素；在后一情形，VI 第 15 页的推论 2 表明 z 大于 x。反之，设该条件成立：若 \(0 \leqslant y \leqslant x\) ，则令 \(x = y + z\) ( \(z \geqslant 0\) )，可得 \(x \leqslant y\) 或 \(x \leqslant z\) ；在第一种情形有 x = y；在第二种情形有 \(x \leqslant x - y\) ，故 \(y \leqslant 0\) ，从而 y = 0；这表明 x 确实是不可约的。

命题 14（DIV）。——要使 A 中非零元素 p 不可约，只需：p 不是单位，且 p 不会在整除 A 中两个元素之积而不整除其中至少一个。当 \(P^{*}\) 是格序的时，此条件也是必要的。

注记。——命题 14（DIV）也可表述如下：若 p 是 A 中非零元素，使理想 \((p)\) 是素理想（I，第 117 页，定义 3），则 p 不可约；反之，若 \(P^{*}\) 是格序的且 p 不可约，则理想 \((p)\) 是素理想。

命题 15.——在序群 G 中，设 \((p_{\iota})_{\iota \in I}\) 是 G 中满足条件 \((P)\) （从而不可约）的两两不同的正元素族。则映射

\[
(n _ {i}) _ {i \in I} \mapsto \sum_ {i \in I} n _ {i} p _ {i}
\]

是序群 \(\mathbf{Z}^{(I)}\) （序群 Z 的直和，VI，第 7 页）到由 \(p_{i}\)

生成的 G 的序子群上的一个同构。只需证明关系式 \(\sum_{\iota \in I} n_{\iota} p_{\iota} \geqslant 0\) 等价于 \(n_{\iota} \geqslant 0\) 对所有 \(\iota\) 成立，

特别地，关系式 \(\sum_{\iota \in I} n_{\iota} p_{\iota} = 0\) 将蕴含 \(n_{\iota} = 0\) 对所有 \(\iota\) 成立，因此这将

表明族 \((p_{\iota})\) 线性无关。现设 \(I'\) （相应地 \(I''\) ）为 \(I\) 的有限子集，由那些 \(\iota\) （使 \(n_{\iota} > 0\) ，相应地使 \(n_{\iota} < 0\) ）组成；我们有

\[
\sum_ {i \in I ^ {\prime}} n _ {i} p _ {i} \geqslant \sum_ {i \in I ^ {\prime \prime}} (- n _ {i}) p _ {i}.
\]

特别地，对 \(\lambda \in I''\) ，这蕴含 \(p_{\lambda} \leqslant \sum_{i \in I'} n_i p_i\) ，并且由归纳法

从性质 (P) 可知，必有 \(p_{\lambda} \leqslant p_{\iota}\) 对某个 \(\iota \in I'\) 成立；由于 \(p_{\iota}\) 不可约，这将蕴含 \(p_{\lambda} = p_{\iota}\) ，这是荒谬的。因此 \(I''\) 为空，这就证明了命题。

定理 2. --- 设 G 是一个滤过群。则以下性质等价：

\begin{enumerate}
\def\labelenumi{\alph{enumi})}
\item
  G 同构于形如 \(\mathbf{Z}^{(I)}\) 的序群。
\item
  G 是格序的，并满足以下条件：
\end{enumerate}

（MIN）\(\mathbf{G}\) 中由正元素组成的每个非空集合都有一个极小元素。

\begin{enumerate}
\def\labelenumi{\alph{enumi})}
\setcounter{enumi}{2}
\item
  G 满足条件 (MIN)，且 G 的每个不可约元都具有性质 (P)。
\item
  G 由其不可约元生成，且 G 的每个不可约元都具有性质 (P)。
\end{enumerate}

首先证明 \(a\) 蕴含 \(b\) 。群 \(\mathbf{Z}^{(I)}\) 是格序的，因为它是全序群的直和。另一方面，设 E 为 \(\mathbf{Z}^{(I)}\) 中正元素组成的一个非空集合，且令 \(x = \sum_{i} n_i e_i\) 为 E 的一个元素（其中

\((e_{i})\) 表示 \(\mathbf{Z}^{(I)}\) 的自然基）；只有有限个（数目为 \(\prod_{i}(n_{i} + 1)\) ）

元素 \(y\)（在 \(\mathbf{Z}^{(I)}\) 中）满足 \(0 \leqslant y \leqslant x\) ，因此集合 \(F\) （由 \(E\) 中小于或等于 \(x\) 的元素组成）更是有限的；由于它非空，它包含一个极小元素（集合论，III，第 170 页，推论 2），这显然是 \(E\) 的一个极小元素。

显然 \(b)\) 蕴含 \(c)\) ，由命题 14。让我们证明 \(c)\) 蕴含 \(d)\) 。由于 \(G\) 是滤过的，只需（VI，第 4 页，命题 4）验证：等于不可约元之和的正元素所成的集合 \(F\) （\(G\) 的正元素之集的子集）等于 \(G_{+} - \{0\}\) 。若非如此，则由 (MIN)，\(F\) 在 \(G_{+} - \{0\}\) 中的补集将有一个极小元素 \(a\) ；按定义 \(a\) 不是不可约的，故而是两个严格正元素 \(x\) 与 \(y\) 之和；由于 \(x < a\) 且 \(y < a\) ，这些元素属于 \(F\) ，因而是不可约元之和，由此可知 \(a\) 也是，这与假设矛盾。最后，\(d)\) 蕴含 \(a)\) ，由命题 15。

我们将把定理 2 应用于主理想整环（VII，第 4 页）和唯一分解整环（AC，VII，第 3 节）中的整除性理论，以及戴德金环（AC，VII，第 2 节）中理想的研究。

\section{序域}

\subsection{序环}

定义 1.——给定交换环 A，我们称 A 上的一个序与 A 的环结构相容，如果它与 A 的加法群结构相容，并且满足以下公理：

(OR) 关系 \(x \geqslant 0\) 与 \(y \geqslant 0\) 蕴含 \(xy \geqslant 0\) 。

环 A 连同这样一个序称为序环。

例子。—— 1）环 Q 和 Z 赋予通常的序，是序环。2）序环的乘积，赋予积序，是序环。特别地，从集合 E 到序环 A 的映射全体所成的环 \(A^{E}\) 是一个序环。

\begin{enumerate}
\def\labelenumi{\arabic{enumi})}
\setcounter{enumi}{2}
\tightlist
\item
  序环的子环，赋予诱导序，也是序环。
\end{enumerate}

在序环中，关系 \(x \geqslant y\) 与 \(z \geqslant 0\) 蕴含 \(xz \geqslant yz\) 。事实上，这些不等式分别等价于 \(x - y \geqslant 0\) , \(z \geqslant 0\) 与 \((x - y)z \geqslant 0\) 。

类似地，可以证明关系 \(x \leqslant 0\) 与 \(y \geqslant 0\) （相应地 \(y \leqslant 0\) ）蕴含 \(xy \leqslant 0\) （相应地 \(xy \geqslant 0\) ）。这些结果常以符号法则的名义被引用（两个元素若同为 \(\geqslant 0\) 或同为 \(\leqslant 0\) ，则称它们同号）。由此推出，若 A 是一个全序环，则每个平方都是正的，特别地，每个幂等元（例如单位元）都是正的。

例。—— 在 Z 上只存在一种全序环结构：事实上 1 \textgreater{} 0，从而 n \textgreater{} 0 对每个自然数 \(n \neq 0\) 成立，由归纳法可得。相形之下，Z 上存在并非全序的序环结构（见下文）。

设 P 为序环 A 的正元素集。已知（VI，第 3 页，命题 3）P 决定 A 上的序关系。说 A 是一个序环，等价于说 P 满足以下性质：

\[
\begin{array}{l l} \left(\mathrm{AP} _ {\mathrm{I}}\right) & \mathrm{P} + \mathrm{P} \subset \mathrm{P} \\ \left(\mathrm{AP} _ {\mathrm{II}}\right) & \mathrm{PP} \subset \mathrm{P} \\ \left(\mathrm{AP} _ {\mathrm{III}}\right) & \mathrm{P} \cap (- \mathrm{P}) = \{0 \}. \end{array}
\]

确实，\((\mathrm{AP}_{\mathrm{I}})\) 与 \((\mathrm{AP}_{\mathrm{III}})\) 断言 A 的加法群是一个序群（VI，第 3 页，命题 3），而 \((\mathrm{AP}_{\mathrm{II}})\) 是 (OR) 的翻译。

回忆一下，A 上的序关系为全序的必要且充分条件是：

\[
\left(\mathrm{AP} _ {\mathrm{IV}}\right) \quad \mathrm{P} \cup (- \mathrm{P}) = \mathrm{A}.
\]

例. --- 在 Z 中，若取 P 为（通常意义下的）正偶数集合，则得到一个并非全序的环。

还要记得，在全序阿贝尔群中，关系 \(n \cdot x = 0\) （对自然数 \(n \neq 0\) ）蕴含 \(x = 0\) （VI，第 4 页）；由此我们得到以下结果。

命题 1. --- 全序环作为 Z-模是无挠的（II，第 313 页）。

\subsection{序域}

定义 2.——交换域配备一个全序，若该序与其环结构相容，则称为序域。

我们仅考虑域上的全序关系，因为其他的序关系非常“病态”（VI，第 38 页，习题 6）。

例子。—— 1）有理数域 \(\mathbf{Q}\) 是一个序域。

\begin{enumerate}
\def\labelenumi{\arabic{enumi})}
\setcounter{enumi}{1}
\item
  序域的子域，赋予诱导序，也是一个序域。
\item
  * 实数域是一个序域。 *
\end{enumerate}

设 \(\mathbf{K}\) 是一个序域。对所有 \(x\in \mathbf{K}\) 我们令

\[
\begin{array}{l l} \operatorname{sgn} (x) = 1 & \text {if} \quad x > 0, \\ \operatorname{sgn} (x) = - 1 & \text {if} \quad x <   0, \\ \operatorname{sgn} (x) = 0 & \text {if} \quad x = 0. \end{array}
\]

那么我们有 \(\operatorname{sgn}(xy) = \operatorname{sgn}(x)\operatorname{sgn}(y)\) ；我们称 \(\operatorname{sgn}(x)\) 为 \(x\) 的符号。映射 \(x \mapsto \operatorname{sgn}(x)\) 从 \(\mathbf{K}^*\) 到乘法群 \(\{-1, +1\}\) 是一个满射同态，其核，即 \(\mathbf{K}\) 的严格正元素所成之集，是 \(\mathbf{K}^*\) 的一个指数为 2 的子群。

反之，若 K 是一个交换域，且 \(s: K^{*} \to \{-1, +1\}\) 是一个核在加法下封闭的满同态，则 s 是某个唯一的序域结构的符号映射，其中严格正元素之集就是 s 的核。

对所有 \(x\) 与 \(y\) 属于 \(\mathbf{K}\) ，我们有 \(x = \operatorname{sgn}(x)|x|\) 且 \(|xy| = |x||y|\) 。

另一方面，每个序域的特征都为零（命题 1）。

命题 2. --- 设 A 是一个全序整环，K 是它的分式域。则 K 上存在唯一一个序，它在 A 上限制为给定的序，并使 K 成为一个序域。

每个 \(x \in K\) 都可表为 \(x = ab^{-1}\) ，其中 a、b 属于 A 且 \(b \neq 0\) 。若 x 为正，则 a 与 b 同号，反之亦然。由此看出，若 K 上存在满足所述条件的序，则它是唯一的，且正元素之集 P 恰为所有 \(ab^{-1}\) 的集合，其中 a、b 是 A 中同号的元素且 \(b \neq 0\) 。剩下要证明 P 满足条件 \((\mathrm{AP}_{\mathrm{I}}), (\mathrm{AP}_{\mathrm{II}}), (\mathrm{AP}_{\mathrm{III}})\) 与 \((\mathrm{AP}_{\mathrm{IV}})\) 。这对 \((\mathrm{AP}_{\mathrm{II}})\) 与 \((\mathrm{AP}_{\mathrm{IV}})\) 是显然的。对 \((\mathrm{AP}_{\mathrm{I}})\) ，考虑 \(ab^{-1} + cd^{-1}\) ，其中可设 \(a, b, c\) 与 \(d\) 为正；这个和是 \((ad + bc)(bd)^{-1}\) ，且 \(ad + bc\) 与 \(bd\) 为正。

为证 \((\mathrm{AP}_{\mathrm{III}})\) ，考虑一个形如 \(ab^{-1} = -cd^{-1}\) 的恒等式，从而 \(ad + bc = 0\) 。若设 a 与 b 同号且 c 与 d 同号，则符号法则表明 ad 与 bc 同号；于是 ad = bc = 0，故 a = c = 0；从而 P 确实满足 \((\mathrm{AP}_{\mathrm{III}})\) 。

例。--- 由于 Z 只容许一个全序环结构（VI，第 19 页，例），域 Q 只容许一个使它成为序域的序：这就是通常的序。

\subsection{序域的扩张}

定义 3.——设 K 是一个序域。K 的一个序扩张是一个二元组 (E, u)，其中 E 是一个序域，u 是从 K 到 E 的一个增同态。

设 K 为一个域，E 为一个序域，且设 \(u: K \rightarrow E\) 是一个同态。关系

\[
x \leqslant y \quad \text { if } \quad u (x) \leqslant u (y)
\]

是 K 上的一个全序关系，使 K 具有一个序域结构，称该结构由 E 的序诱导。若 K 与 E 都是序域，则同态 \(u: K \rightarrow E\) 是增同态，当且仅当 K 的序域结构由 E 的序域结构诱导。我们通常把 K 与它在 u 下的像等同起来。

例子. --- 1) 每个序域 K 都是 Q 的一个序扩张。事实上，K 是 Q 的扩张，因为其特征为零；另一方面，正如我们刚刚看到的，Q 只能以一种方式定序。

\begin{enumerate}
\def\labelenumi{\arabic{enumi})}
\setcounter{enumi}{1}
\tightlist
\item
  设 K 是一个序域，并设 K(X) 为 K 上关于单个未定元的有理函数域。让我们在多项式环 K{[}X{]} 上定义一个序：取正元素为 0 以及首项系数为正的那些多项式。这样我们得到一个全序环，其序延拓了 K 上的序。应用命题 2，我们赋予 K(X) 一个 K 的序扩张的结构。* 对 K = R，可以证明如此定义在 K(X) 上的序关系就是在 +∞ 附近的增长序（参见 VI，第 24 页，命题 4）。*
\end{enumerate}

定理 1. --- 对于 K 的扩张 E，要使 E 能够具有 K 的序扩张的结构，以下条件是必要且充分的：

（OE）关系 \(p_{1}x_{1}^{2} + \cdots + p_{n}x_{n}^{2} = 0\) 蕴含

\[
p _ {1} x _ {1} = \dots = p _ {n} x _ {n} = 0
\]

对元素对 \((x_{i},p_{i})\) 的任意有限序列，其中 \(x_{i}\) 属于 E，\(p_i\) 是 K 的正元素。

条件 (OE) 显然等价于：

(OE') 元素 \(-1\) 不是形如 \(px^2\) 的元素之和 ( \(x \in \mathbf{E}\) , \(p \in \mathbf{K}\) , \(p \geqslant 0\) )。

条件 (OE) 是必要的：如果 E 是 K 的一个序扩张，则元素 \(p_{i}x_{i}^{2}\) 在 E 中是正的，故当它们的和为零时它们为零。另一方面，\(p_{i}x_{i}^{2}=0\) 等价于 \(p_{i}x_{i}=0\) 。

反之，假设条件 (OE) 成立，则我们将通过在 E 中构造一个满足条件 \((\mathrm{AP}_{\mathrm{I}})\) , \((\mathrm{AP}_{\mathrm{II}})\) , \((\mathrm{AP}_{\mathrm{III}})\) 与 \((\mathrm{AP}_{\mathrm{IV}})\) 且包含 K 的正元素之集 \(K_{+}\) 的子集 P 来在 E 上定义一个序。这样的子集 P 必使 E 成为 K 的一个序扩张，因为我们会有 \(K \cap P = K_{+}\) ；事实上，如果 P 包含 K 的一个元素 -a \textless{} 0，则 a 将属于 \(P \cap (-P)\) ，这与 \((\mathrm{AP}_{\mathrm{III}})\)

相矛盾。为了定义 P，让我们考虑由 E 中满足 \((\mathrm{AP}_{\mathrm{I}})\) , \((\mathrm{AP}_{\mathrm{II}})\) 与 \((\mathrm{AP}_{\mathrm{III}})\) 且包含 \(K_{+}\) 与 E 的平方元素所成之集 C 的并的子集构成的集合 M。这个集合 M 是非空的，因为它包含集合 \(P_{0}\) ，即形如 \(\sum_{i} p_{i} x_{i}^{2}\) 的元素所成之集（\(P_{0}\) 满足 \((\mathrm{AP}_{\mathrm{III}})\) 这一事实由

（OE）立即可得）。

此外，M 是归纳的（集合论，III，第 154 页，定义 3）。因此，根据集合论，III，第 154 页的定理 2，M 中存在一个极大元，我们尚需证明它满足 \((\mathrm{AP}_{\mathrm{IV}})\) ；而这由以下引理得出：

引理。--- 设 \(P \in M\) 且 \(x \notin P\) ；则存在 \(P' \in M\) 使得 \(P \subset P'\) 且 \(-x \in P'\) 。

取 \(P' = P - xP\) ，并验证 \(P'\) 具有所需的性质。由于 \(0 \in C \subset P\) ，我们有 \(P \subset P'\) 。由此 \(C \subset P'\) 且 \(K_{+} \subset P'\) 。由于 \(1 \in C \subset P\) ，我们有 \(-x \in P'\) 。我们有

\[
\mathrm{P} ^ {\prime} + \mathrm{P} ^ {\prime} = \mathrm{P} - x \mathrm{P} + \mathrm{P} - x \mathrm{P} = \mathrm{P} + \mathrm{P} - x (\mathrm{P} + \mathrm{P}) \subset \mathrm{P} - x \mathrm{P} = \mathrm{P} ^ {\prime},
\]

由此得 \((\mathbf{AP}_1)\) 。我们有

\[
\begin{array}{r l} \mathrm {P^ {\prime} P^ {\prime}} = & (\mathrm{P} - x \mathrm{P}) (\mathrm{P} - x \mathrm{P}) \subset \\ & \subset \mathrm{PP} + x ^ {2} \mathrm{PP} - x (\mathrm{PP} + \mathrm{PP}) \subset \mathrm{P} + \mathrm{CP} - x \mathrm{P} \subset \mathrm{P} - x \mathrm{P} = \mathrm {P^ {\prime}}, \end{array}
\]

由此得 \((\mathrm{AP}_{\mathrm{II}})\) 。最后，验证 \((\mathrm{AP}_{\mathrm{III}})\) ：设给定一个形如 \(p - xq = -(r - xs)\) 的恒等式，其中 p、q、r、s 属于 P；由此我们推出关系 \(x(s + q) = p + r\) ；若 \(s + q \neq 0\) ，我们有

\[
x = (s + q) ^ {- 2} (s + q) (p + r) \in \mathrm{CPP} \subset \mathrm{P},
\]

与假设相反；因此 \(s + q = 0\) ，从而 \(p + r = 0\) ；由于 \(\mathbf{P}\) 满足 \((\mathbf{AP}_{\mathrm{III}})\) ，我们推断出 \(s = q = r = p = 0\) ，这就完成了证明。

推论 1（“阿廷–施赖尔定理”）。——在交换域 E 上存在一个序关系使其成为序域的充分必要条件是：关系 \(x_{1}^{2} + \cdots + x_{n}^{2} = 0\) 蕴含 \(x_{1} = \cdots = x_{n} = 0\) 。

必要性是显然的。反过来，所述条件蕴含 E 的特征为零，因此 E 是 Q 的扩张；于是条件 (OE) 得到满足，且定理 1 表明在 E 上存在 Q 的序扩张结构，即一个序域结构。

在 -1 为平方数的域 E 上不存在任何序域结构，特别地，在代数闭域上也是如此。

推论 2. --- 设 E 是 K 的一个扩张，且 E 容许具有 K 的序扩张的结构。元素 \(x \in E\) 在 E 上的每一个这样的结构下为正，当且仅当 x 具有形式 \(\sum_{i} p_{i} x_{i}^{2}\) ，其中 \(x_{i} \in E\) ，而

\(p_i\) 是 \(\mathbf{K}\) 的正元素。

该条件显然充分；它也是必要的，因为（采用定理 1 证明中的记号），如果 \(x \notin \mathbf{P}_0\) ，则存在一个极大元 \(\mathbf{P}\) 属于 \(\mathfrak{M}\) ，使得 \(x \notin \mathbf{P}\) ；于是根据引理，\(-x \in \mathbf{P}\) ，且 \(x\) 在由 \(\mathbf{P}\) 定义的序下不为正，因为 \(x \neq 0\) 。

\subsection{序域的代数扩张}

设 K 为一个序域，且 f 为 K{[}X{]} 中的多项式。我们说 f 在 K 中变号，如果存在 K 中的两个元素 a 和 b，使得 \(f(a)\) \(f(b) < 0\) ；这时我们就说 f 在 a 与 b 之间变号。

命题 3.——设 K 是一个序域，f 是 K 上的不可约多项式，且 f 在 K 中的 a 与 b 之间变号。则扩张 \(E = K[X]/(f)\) 容许一个 K 的序扩张的结构。

我们通过对 f 的次数 n 进行归纳来论证。当 n = 1 时，证明是平凡的。假设结论对次数 \(\leqslant n - 1\) 的情形成立，我们用反证法证明次数为 n 的情形；由定理 1，我们因此假设一个形如

\(1 + \sum_{i}p_{i}f_{i}^{2}(\mathbf{X})\equiv 0 (\mathrm{mod} f(\mathbf{X}))\) ，其中 \(f_{i}\in \mathbf{K}[\mathbf{X}], p_{i}\in \mathbf{K}\) 且 \(p_i\geqslant 0\)

不失一般性，我们可以假设这些 \(f_{i}\) 的次数 \(\leqslant n - 1\) （IV，第 11 页，推论）。于是

\[
1 + \sum_ {i} p _ {i} f _ {i} ^ {2} (\mathbf {X}) = h (\mathbf {X}) f (\mathbf {X})
\]

，其中 \(h \neq 0\) 的次数至多为 n - 2。在上面的不等式中把 X 换成 a 和 b，我们看到 \(h(a) f(a) > 0\) 且 \(h(b) f(b) > 0\) 。由于根据假设 f 在 a 与 b 之间变号，我们得出 \(h(a) h(b) < 0\) 。于是对不可约因子 \(g(\mathbf{X})\) （\(h(\mathbf{X})\) 的因子之一）有类似的不等式：即 \(g(a)g(b) < 0\) 。但 \(1 + \sum_{i}p_{i}f_{i}^{2}(\mathbf{X})\equiv 0 (\mathrm{mod} g(\mathbf{X}))\) ，这表明域 \(\mathrm{K}[\mathrm{X}] / (g)\) 不能是 \(\mathbf{K}\) 的序扩张（定理 1），与归纳假设相矛盾。

注记。--- 存在序域 K 上的不可约多项式 f，它们在 K 中不变号，但 \(\mathrm{K}[\mathrm{X}]/(f)\) 却容许具有 K 的序扩张的结构（参见 VI，第 43 页，习题 26，c））。

为了应用前面的命题，我们需要以下结果：

命题 4. --- 设 K 是一个序域，并设 \(f \in \mathbf{K}[\mathbf{X}]\) 。在 K 中存在一个区间，在该区间的补集上， \(f\) 与其最高次项具有相同的符号。

我们可以立即将问题化归为多项式为首一多项式的情形；于是可以写出 \(f(x) = x^{n}(1 + a_{1}x^{-1} + \cdots + a_{n}x^{-n})\) ，对 \(x \neq 0\) 。设

\[
\mathbf {M} = \sup \left(1, | a _ {1} | + \dots + | a _ {n} |\right).
\]

对于 \(|x| > M\) 我们有 \(1 + a_{1}x^{-1} + \cdots + a_{n}x^{-n} > 0\) ，这就完成了命题的证明。

推论 1. --- 序域的奇数有限次扩张的每一个扩张都容许序扩张的结构。

这样的扩张，由于是单生成的（V，第 40 页，定理 1），同构于 K{[}X{]}/(f)，其中 f 是一个奇数次的不可约多项式。那么只需证明 f 在 K 中变号（命题 3），这由命题 4 立即得出。

推论 2. —— 若 a 是序域 K 的正元素，则多项式 \(X^{2}-a\) 的每个分裂域 E 都容许具有 K 的序扩张的结构。

若 a 是 K 中的平方，则结论是平凡的。否则，多项式 \(f(\mathbf{X}) = \mathbf{X}^{2} - a\) 是不可约的且变号，因为 \(f(0) < 0\) ，且对 K 的某个区间的补集中的 x，\(f(x)\) 与 \(x^{2}\) 同号，即为正。我们现在可以通过应用命题 3 来完成证明。

注：--- 当序域 K 包含 K 的正元素 \(a\) 的“平方根”（多项式 \(\mathbf{X}^2 - a\) 的根）时，记号 \(\sqrt{a}\) 通常保留给正的平方根。若 K 不包含 \(b\) 与 \(-b\) （\(a\) 在域 E 中的平方根），则后者可以以两种方式成为 K 的序扩张，每种方式都由另一个通过把 \(b\) 映为 \(-b\) 的 K-自同构诱导而来。选择这两个序中的哪一个，便确定了 \(\sqrt{a}\) ：它是元素 \(b\) 与 \(-b\) 中为正的那一个。

如果 \(\pmb{a}\) 与 \(a'\) 是 \(\mathbf{K}\) 的两个正元素，其平方根位于 \(\mathbf{K}\) ，则 \(\sqrt{aa'} = \sqrt{a}\sqrt{a'}\) ，这由 \(\sqrt{a}\) 的定义以及符号法则得出。

\subsection{极大序域}

定义 4.——序域 K 是极大的，如果 K 的每一个序代数扩张都是平凡的。

例。--- * 我们稍后将看到（一般拓扑学，VIII，第 1 页）实数域 R 是一个极大序域。 *

极大序域的存在性是以下定理的一个推论：

定理 2.——每个序域 K 都容许一个序代数扩张，该扩张是一个极大序域。

可以证明，这个序扩张在 K-同构意义下是唯一的（VI，第 40 页，习题 15）。

设 \(\Omega\) 为 \(K\) 的一个代数闭包，并设 \(\mathfrak{N}\) 为由二元组 \((A, \omega)\) 组成的集合，其中 \(A\) 是 \(K\) 的一个扩域，包含于 \(\Omega\) ，而 \(\omega\) 是 \(A\) 上的一个序，它使 \(A\) 成为 \(K\) 的一个序扩张。给 \(\mathfrak{N}\) 定序，所用关系为 « \(L\) 是 \(M\) 的序扩张» （介于 \(M\) 与 \(L\) 之间）。配备此序后， \(\mathfrak{N}\) 是一个归纳有序集：事实上，若 \((L_{t})\) 是 \(\mathfrak{N}\) 中元素的一个全序族，则域 \(L = \cup_{t} L_{t}\) ，取 \(L_{+} = \cup_{t} (L_{t})_{+}\) 定序后，是

\(L_{i}\) 的一个上界。于是根据集合论，III，第 154 页，定理 2，N 有一个极大元，这就完成了证明。

命题 5. --- 设 K 是一个极大序域，并设 f 是 K{[}X{]} 中的多项式，它在 K 的两个元素 a 和 b（其中 a \textless{} b）之间变号。那么 f 在 K 中有一个根 x，使得 a \textless{} x \textless{} b。

f 的至少一个不可约因子，设为 h，在 a 与 b 之间变号。于是域 \(\mathrm{K}[\mathrm{X}]/(h)\) 容许 K 的一个序扩张的结构（VI，第 23 页，命题 3），且 h 的次数为 1（定义 4）。由于 \(h(a)h(b)<0\) ，h 的唯一根 x 满足 a\textless x\textless b，因为一次多项式函数是单调的。

命题 6. --- 极大序域 K 的每个正元素在 K 中都有平方根。K{[}X{]} 中每个奇数次多项式在 K 中至少有一个根。

这直接由 VI 第 24 页命题 4 的推论 2 和 1 得出。

推论。--- 在极大序域 K 上，只存在一种与域结构相容的序。

确实，\(\mathbf{K}\) 的正元素由其代数结构决定：它们就是平方数。

\subsection{极大序域的特征刻画。欧拉-拉格朗日定理}

VI 第 25 页命题 6 所表达的性质刻画了极大序域。更精确地说：

定理 3（欧拉-拉格朗日）。--- 设 K 为一个序域。则以下三个性质等价：

\begin{enumerate}
\def\labelenumi{\alph{enumi})}
\item
  域 \(\mathbf{K}(i)\) 是代数闭的（其中 \(i\) 表示 \(-1\) 的一个平方根）。
\item
  序域 \(\mathbf{K}\) 是极大的。
\item
  \(\mathbf{K}\) 的每个正元素都是平方，且 \(\mathbf{K}[\mathbf{X}]\) 的每个奇数次多项式在 \(\mathbf{K}\) 中有一个根。
\end{enumerate}

显然 \(a)\) 蕴含 \(b)\) ：事实上，\(K\) 在同构意义下只有两个代数扩张，即域 \(K\) 其自身及 \(K(i)\) ，后者无法定序，因为 \(-1\) 是一个平方。

\(b)\) 蕴含 \(c)\) 这一事实无非就是 VI 第 25 页的命题 6。

我们尚需证明 \(c)\) 蕴含 \(a)\) 。这将由接下来的两个命题得出。

命题 7. --- 设 K 是一个序域，其中每个正元素都是平方。则 K(i) 的每个元素都是平方，且 K(i) 上每个次数为 2 的多项式在 K(i) 中都有根。

让我们首先证明第二个断言可以化归为第一个。可以把二次多项式 \(aX^{2} + bX + c\) ( \(a \neq 0\) )写成如下形式，通常称为标准三项式形式：

\[
a \left(\left(X + (b / 2 a)\right) ^ {2} - \left(b ^ {2} - 4 a c\right) / 4 a ^ {2}\right).
\]

如果 \(d\) 是 \((b^2 - 4ac) / 4a^2\) 的一个平方根，则 \(d - (b / 2a)\) 是所考虑的二次多项式的根。

现在我们证明每个元素 \(a + bi\) ( \(a \in \mathbf{K}\) , \(b \in \mathbf{K}\) ) 都是平方；我们要寻找一个元素 \(x + yi\) 使得

\[
(x + y i) ^ {2} = a + b i;
\]

这转化为 \(x^{2} - y^{2} = a\) 与 \(2xy = b\) ，由此我们推断出

\[
(x ^ {2} + y ^ {2}) ^ {2} = a ^ {2} + b ^ {2}.
\]

设 \(c\) 表示 \(a^2 + b^2\) 的正平方根；则 \(c \geqslant |a|\) , \(c \geqslant |b|\) 且 \(x^2 + y^2 = c\) ，由此得出 \(x^2 = (c + a)/2\) 与 \(y^2 = (c - a)/2\) 。由于 \(c \geqslant |a|\) ，这些方程在 K 中可解，且若 \(x_0\) 与 \(y_0\) 是两个解，则 \(x_0^2 - y_0^2 = a\) 且 \(2x_0y_0 = \pm b\) ，我们通过取 \(x = x_0\) 与 \(y = b/2x_0\) 得到所需的平方根。

命题 8. --- 设 K 是一个交换域（特征任意），并设 \(\mathbf{K}'\) 是多项式 \(\mathbf{X}^2 + 1 \in \mathbf{K}[\mathbf{X}]\) 的一个分裂域（V，第 21 页）。假设：

\begin{enumerate}
\def\labelenumi{\alph{enumi})}
\item
  K{[}X{]} 中每个奇数次的多项式在 K' 中有根；
\item
  \(\mathbf{K}'[\mathbf{X}]\) 中每个次数为 2 的多项式在 \(\mathbf{K}'\) 中有根。则 \(\mathbf{K}'\) 是代数闭的。
\end{enumerate}

首先注意，只需证明 K{[}X{]} 中每个非常数多项式在 K' 中有根：若 \(K' = K\) 这是清楚的；若 \(K' \neq K\) ，则 \([K': K] = 2\) ；设 \(a \mapsto \bar{a}\) 表示 \(K'\) 的唯一异于恒等映射的 K-自同构；若 \(f \in K'[X]\) 且 \(\bar{f}\) 表示对 f 的系数应用 \(a \mapsto \bar{a}\) 所得到的多项式，则 \(f\bar{f} \in K[X]\) ；若 \(a \in K'\) 是 \(f\bar{f}\) 的根，则 a 或是 f 的根，或是 \(\bar{f}\) 的根；因此 a 或 \(\bar{a}\) 是 f 的根。

因此设 f 是 K 上的一个多项式，次数为 \(2^{n}p\) ，p 为奇数。我们将对 n 进行归纳，该性质在 n = 0 时由假设 a) 成立。设 E 是 K 的一个扩域，在其中 f 分裂为线性因子：

\[
f (\mathbf {X}) = \prod_ {i} \left(\mathbf {X} - a _ {i}\right).
\]

设 \(b \in K\) ；令 \(y_{ij} = a_i + a_j + ba_i a_j \in E\) 且

\[
h (\mathbf {X}) = \prod_ {i <   j} \left(\mathbf {X} - y _ {i j}\right) \in \mathrm{E} [ \mathbf {X} ].
\]

该多项式的系数是这些 \(a_i\) 的、系数在 \(\mathbf{K}\) 中的对称函数；因此它属于 \(\mathbf{K}[\mathbf{X}]\) （IV，第 62 页，定理 1）；由于它的次数为 \(2^n p(2^n p - 1)/2 = 2^{n-1}p'\) ( \(p'\) 奇数)，根据归纳假设它有一个根 \(y_{ij}\) 属于 \(\mathbf{K}'\) 。如果我们注意到这对所有 \(b \in \mathbf{K}\) 成立，并且 \(\mathbf{K}\) 是无限域（事实上，一个有限域，由于它具有任意大奇数次数的单生成扩张（V，第 94 页，命题 3），不能满足 \(a\) )），那么我们可以推断出至少存在一对 \((i,j)\) 使得

\[
a _ {i} + a _ {j} + b a _ {i} a _ {j} \in \mathbf {K} ^ {\prime} \quad \text { and } \quad a _ {i} + a _ {j} + b ^ {\prime} a _ {i} a _ {j} \in \mathbf {K} ^ {\prime},
\]

，其中 \(b \neq b'\) 。则 \(a_i + a_j\) 与 \(a_i a_j\) 是 \(K'\) 的元素，因此 \(a_i\) 与 \(a_j\) 也是，因为它们是二次方程

\[
x ^ {2} - \left(a _ {i} + a _ {j}\right) x + a _ {i} a _ {j} = 0.\tag{Q.E.D.}
\]

的根。关于命题 8 的推广以及基于伽罗瓦理论的另一种证明，参见 VI，第 46 页，习题 33。

设 K 是一个序域，并设 \(\mathbf{K}^{\prime} = \mathbf{K}(i)\) ；对于每个元素 \(z = a + bi\) 属于 \(K^{\prime}\) ，范数 \(z\overline{z} = a^{2} + b^{2}\) （相对于 K，III，第 544 页，例 1）是 K 的一个正元素，且仅当 z = 0 时为零。如果 K 中的每个正元素都是平方（特别地，若 K 是极大序域），则范数 \(z\overline{z}\) 的正平方根称为 z 的绝对值，记作 \(|z|\) 。由于 \(|zz^{\prime}|^{2} = |z|^{2} |z^{\prime}|^{2}\) ，我们有 \(|zz^{\prime}| = |z| \cdot |z^{\prime}|\) 。

此外，三角不等式

\[
\left| z + z ^ {\prime} \right| \leqslant \left| z \right| + \left| z ^ {\prime} \right|
\]

对每一对元素 \(z, z'\) 属于 \(\mathbf{K}'\) 都成立。事实上，如果 \(z = a + bi\) 且 \(z' = a' + b'i\) ，那么这个不等式等价于

\[
(a + a ^ {\prime}) ^ {2} + (b + b ^ {\prime}) ^ {2} \leqslant a ^ {2} + b ^ {2} + a ^ {\prime 2} + b ^ {\prime 2} + 2 \sqrt {(a ^ {2} + b ^ {2}) (a ^ {\prime 2} + b ^ {\prime 2})}
\]

因而也等价于

\[
(a a ^ {\prime} + b b ^ {\prime}) ^ {2} \leqslant (a ^ {2} + b ^ {2}) (a ^ {\prime 2} + b ^ {\prime 2})
\]

它可以写成 \((ab' - ba')^2 \geqslant 0\) 。

定理 3 使我们能够确定一个极大序域上的所有不可约多项式：

命题 9. —— 若 K 是一个极大序域，则 K{[}X{]} 中仅有的不可约多项式是一次多项式，以及二次多项式 \(aX^{2} + bX + c\) ，其中 \(b^{2} - 4ac < 0\) 。

由于 \(\mathbf{K}(i)\) 是代数闭的，K 的每个代数扩张、因而 K 上的每个不可约多项式，其次数为 1 或 2。为了看出哪些二次多项式是不可约的，只需考虑标准形式 \(a((\mathbf{X} + (b/2a))^2 - (b^2 - 4ac)/4a^2)\) （参见 VI，第 26 页，命题 7）。

注：--- 转化为标准三项式形式可得到这一更强的结果：多项式 \(aX^{2} + bX + c\) 在给定的序域 K 上具有常号，当且仅当 \(b^{2} - 4ac < 0\) ，且此时多项式的符号与 a 的符号相同。

\subsection{序域上的向量空间}

设 K 为一个序域，E 为 K 上的一个向量空间。关系«存在 \(\lambda > 0\) 在 K 中使得 \(y = \lambda x \gg\) ，在集合 \(E - \{0\}\) 的两个元素 x 与 y 之间，是一个等价关系。在此关系下的等价类称为以 0 为原点的开半直线；一个开半直线与 \(\{0\}\) 的并称为以 0 为原点的闭半直线（有时简称半直线）。每个向量 \(a \neq 0\) ，若它包含在一条开（相应地，闭）半直线 \(\Delta\) 中，则称为 \(\Delta\) 的方向向量，而 \(\Delta\) 是向量 \(\lambda a\) （对所有标量 \(\lambda > 0\) ，相应地 \(\lambda \geqslant 0\) ）所成的集合。每条过 0 的直线 D 恰好包含两条以 0 为原点的开（相应地，闭）半直线；如果 \(\Delta\) 是其中之一，则 \(-\Delta\) 是另一个（称为 \(\Delta\) 的反向半直线）。

现在如果 F 是 K 上的一个仿射空间，且 E 是 F 的平移空间，那么 F 的任何形如 \(\Delta = a + \Delta_{0}\) 的子集，其中 \(\Delta_{0}\) 是 E 的一条开（相应地，闭）半直线，称为以原点 \(a \in F\) 为起点的开（相应地，闭）半直线。半直线 \(\Delta_{0}\) 完全由 \(\Delta\) 确定（因为它是对任意 \(b \neq a\) 属于 \(\Delta\) 而言方向向量为 b - a 的那条半直线），称为 \(\Delta\) 的方向；\(\Delta_{0}\) 的方向向量也称为 \(\Delta\) 的方向向量。

现在假设 E 在 K 上具有有限维数 n；那么已知（III，第 518 页，推论 1）第 n 次外幂 \(\bigwedge^{n}E\) 是 K 上的一维向量空间，因此是以 0 为原点的两条相反闭半直线之并。这些半直线称为 E 的定向；空间 E 连同给定的其中一条半直线 \(\Delta\) 称为已定向的；一个 \(n\) -向量 \(z\) 在此定向下，若它属于 \(\Delta\) （相应地，属于 \(-\Delta\) ），则称为正的（相应地，负的）；在相反的定向下，它是负的（相应地，正的）。

仿射空间 F 在 K 上的一个定向，按定义是 F 的平移空间的一个定向；空间 F 连同这样的定向，称为有向仿射空间。

设 \(\mathbf{E}\) 是 \(\mathbf{K}\) 上一个已定向的向量空间，维数为 \(n\) ；有序基 \((a_i)_{1 \leqslant i \leqslant n}\) （\(\mathbf{E}\) 的）称为正的或顺向的（相应地，负的或逆向的），如果 \(n\) -向量 \(a_1 \wedge a_2 \wedge \ldots \wedge a_n\) 为正（相应地，为负）。若 \(u\) 是向量空间 \(\mathbf{E}\) 的一个自同构，则 \((\bigwedge^{n} u)(z) = \det(u) \cdot z\) 对所有 \(z \in \bigwedge^{n} \mathbf{E}\) 成立，故

\(\bigwedge^{n} u\) 保持 E 的定向不变（或者，如我们所说，保持定向）的一个充分必要条件是 \(\det(u) > 0\) ；具有此性质的自同构恰好是 E 作为已定向向量空间的自同构；它们构成一个正规子群 \(\mathbf{GL}^{+}(\mathbf{E})\) ，即线性群 \(\mathbf{GL}(\mathbf{E})\) 的子群，只要 \(E \neq 0\) ，其指数为 2。

当 \(\mathbf{E} = 0\) 时，\(\mathbf{GL}(\mathbf{E}) = \operatorname{End}(\mathbf{E})\) 仅包含恒等映射 \(1_{\mathbf{E}}\) ，且按定义 \(\det(1_{\mathbf{E}}) = 1\) 。注意，在此情形下按定义 \(\bigwedge^{n}\mathbf{E} = \bigwedge^{0}\mathbf{E} = \mathbf{K}\) ；\(\mathbf{K}\) 的由正标量组成的半直线称为零空间的标准定向。

设 \(\mathbf{M}\) 与 \(\mathbf{N}\) 是维数分别为 \(p\) 与 \(n - p\) 的两个互补子空间，属于向量空间 \(\mathbf{E}\) （维数为 \(n\) ）；如果 \(z'\) （相应地 \(z''\) ）是 \(\bigwedge^{p} \mathbf{M}\) （相应地 \(\bigwedge^{n - p} \mathbf{N}\) ）中的一个非零向量，则 \(z' \wedge z''\) 是

\(\bigwedge^{n} \mathbf{E}\) 中的非零向量。给定 M 上的一个定向和 N 上的一个定向，向量 \(z' \wedge z''\) 对正的 \(z'\) 与 \(z''\) 构成 E 的一个定向，称为 M 的定向与 N 的定向的乘积定向（当 \(p(n - p)\) 为奇数时，它依赖于因子的顺序）。反之，给定 E 和 M 上的定向，存在 N 上唯一的定向，使得 E 上给定的定向是 M 上给定的定向与该 N 上的定向（按此顺序）的乘积；该定向称为关于 E 的定向的 M 的定向的互补定向。若 \(N'\) 是 M 的第二个互补子空间，则平行于 M 的典范投影 \(N \to N'\) 将 N 的互补定向映到 \(N'\) 的互补定向。因此，N 的互补定向在典范映射 \(N \to E/M\) 下的像不依赖于补空间 N 的选择；它称为 E 的定向除以 M 的定向所得的 E/M 上的商定向。

\BourbakiExercises

\BourbakiExerciseGroup

\begin{enumerate}
\def\labelenumi{\arabic{enumi})}
\tightlist
\item
  一个不一定交换的序幺半群 M 是一个（乘法记法写的）幺半群，配备了一个序，使得关系 \(x \leqslant y\) 蕴含 \(zx \leqslant zy\) 与 \(xz \leqslant yz\) 对所有 \(z \in M\) 。设 G 是一个不一定交换的序群。证明：
\end{enumerate}

\begin{enumerate}
\def\labelenumi{\alph{enumi})}
\item
  若 P 表示 G 中大于单位元 e 的元素之集，则 P·P = P， \(P \cap P^{-1} = \{e\}\) 且 \(aPa^{-1} = P\) 对所有 \(a \in G\) 。证明逆命题。找出使 G 成为全序群的一个条件。
\item
  如果元素 \(\sup (x,y)\) , \(\inf (x,y)\) 之一存在，则另一个也存在，并且
\end{enumerate}

\[
\sup (x, y) = x (\inf (x, y)) ^ {- 1} y = y (\inf (x, y)) ^ {- 1} x.
\]

\begin{enumerate}
\def\labelenumi{\alph{enumi})}
\setcounter{enumi}{2}
\tightlist
\item
  元素 \(\sup(x, e)\) 与 \(\sup(x^{-1}, e)\) 交换。两个互素的元素交换。d) 由 G 的不可约元生成的子群 \(G'\) 是交换的。由此推出 VI 第 18 页的定理 2 仍然成立。
\end{enumerate}

\begin{enumerate}
\def\labelenumi{\arabic{enumi})}
\setcounter{enumi}{1}
\tightlist
\item
  设 E 是一个带有合成律 \((x, y) \mapsto xy\) （不一定结合）的格，使得对所有 \(a \in E\) ，映射 \(x \mapsto ax\) 与 \(x \mapsto xa\) 都是 E 的序自同构。设 \(x_{a}\) （相应地 \(_{a}x\) ）表示由 \((x_{a})a = x\) （相应地 \(a(_{a}x) = x\) ）定义的 E 的元素，并假设对每个 \(a \in E\) ，映射 \(x \mapsto a_{x}\) 与 \(x \mapsto_{x} a\) 是从 E 的序到相反序的同构。证明：对任意的 x, y, z，
\end{enumerate}

\[
\left(z _ {\inf (x, y)}\right) \cdot x = \left(z _ {y}\right) \cdot \sup (x, y).
\]

\begin{enumerate}
\def\labelenumi{\arabic{enumi})}
\setcounter{enumi}{2}
\item
  设 \(G\) 是一个序群，其正元素之集 \(P\) 不为 0；证明 \(G\) 是无限群，且不能有最大（或最小）元素。
\item
  设 G 是一个序群，P 为 G 的正元素之集，f 为从 G 到商群 G/H 上的自然映射。要使 \(f(P)\) 使 G/H 成为一个序群，必要且充分条件是 \(0 \leqslant y \leqslant x\) 与 \(x \in H\) 蕴含 \(y \in H\) ；此时 H 称为 G 的凸子群，且 G/H 以此方式被视为一个序群。若 G 是格序的，则 G/H 为格序且 \(f(\sup(x, y)) = \sup(f(x), f(y))\) 成立的必要且充分条件是关系 \(|y| \leqslant |x|\) 与 \(x \in H\) 蕴含 \(y \in H\) ；证明这就是说 H 是一个滤过凸子群。证明一个格序群若除自身与 \(\{0\}\) 外没有其他滤过凸子群，则是全序的（考虑由 G 的两个正元素生成的滤过凸子群）；* 由此推出（一般拓扑学，V，第 25 页，习题 1）G 同构于实数的一个加法子群。*
\item
  给出一个序群的例子，其序是非离散的，且具有非零的有限阶元素（取一个适当的序群关于满足 \(P \cap H = \{0\}\) 的子群 H 的商）。
\item
  设 G 是一个序群，设 P 是其正元素之集。证明 P - P 是 G 的最大滤过子群，并且它是凸的。商上的序关系是什么？
\item
  在群 \(\mathbf{Z}\) 中，如果 \(\mathbf{P}\) 取为由 0 和整数 \(\geqslant 2\) 组成的集合，则所得的序群是滤过的但不是格序的（证明元素 \(x\) （满足 \(x \geqslant 0\) 且 \(x \geqslant 1\) ）所成的集合有两个不同的极小元素）。
\item
  设 \(x\) 是序群中的一个元素，使得 \(y = \inf(x, 0)\) 有定义；如果 \(n > 0\) 是整数，则 \(nx \geqslant 0\) 蕴含 \(x \geqslant 0\) （我们有
\end{enumerate}

\[
n y = \inf (n x, (n - 1) x, \dots , 0) \geqslant \inf ((n - 1) x, \dots , 0) = (n - 1) y);
\]

因此 \(nx = 0\) 蕴含 \(x = 0\) 。

\begin{enumerate}
\def\labelenumi{\arabic{enumi})}
\setcounter{enumi}{8}
\item
  在格序群 G 中，证明一族 \((H_{\alpha})\) 滤过凸子群之和是滤过凸子群（利用 VI，第 12 页，推论）。
\item
  证明，对每个有限序列 \((x_{i})\) ( \(1 \leqslant i \leqslant n\) )（格序群 \(G\) 的元素），
\end{enumerate}

\[
\sup \left(x _ {i}\right) = \sum_ {i} x _ {i} - \sum_ {i <   j} \inf \left(x _ {i}, x _ {j}\right) + \dots + (- 1) ^ {p + 1} \sum_ {i _ {1} <   i _ {2} <   \dots <   i _ {p}} \inf \left(x _ {i _ {1}}, \dots , x _ {i _ {p}}\right) +   + \dots + (- 1) ^ {n + 1} \inf \left(x _ {1}, \dots , x _ {n}\right).
\]

（通过基于命题 7（VI，第 10 页）的归纳论证，并利用上确界对下确界的分配性。）

\begin{enumerate}
\def\labelenumi{\arabic{enumi})}
\setcounter{enumi}{10}
\item
  设 \((x_{i})\) 是格序群 G 的 n 个元素组成的族；对每个整数 k \((0 \leqslant k \leqslant n)\) ，设 \(d_{k}\) （相应地 \(m_{k}\) ）表示 \(\binom{n}{k}\) 个具有互不相同指标的 k 项 \(x_{i}\) 之和的下确界（相应地，上确界）。证明 \(d_{k} + m_{n-k} = x_{1} + x_{2} + \cdots + x_{n}\) 。
\item
  给定格序群 \(\mathbf{G}\) ，一个子群 \(\mathbf{H}\) （\(\mathbf{G}\) 的子群）称为余格，如果对每一对 \(x, y\) （\(\mathbf{H}\) 的元素），有 \(\sup_{\mathbf{G}}(x, y) \in \mathbf{H}\) （从而等于 \(\sup_{\mathbf{H}}(x, y)\) ）。
\end{enumerate}

\begin{enumerate}
\def\labelenumi{\alph{enumi})}
\item
  如果 \(\mathbf{G} = \mathbf{Q} \times \mathbf{Q} \times \mathbf{Q}\) （其中 \(\mathbf{Q}\) 具有通常的序），则子群 \(\mathbf{H}\) （由 \((x, y, z)\) （其中 \(z = x + y\) ）组成）是一个格但不是余格。
\item
  格序群 G 的每个滤过凸（习题 4）子群 H 都是余格。
\item
  设 G 为一个格序群，H 为 G 的任意子群。定义 H' 为 H 的有限子集的下确界之集，H'' 为 H' 的有限子集的上确界之集。证明 H'' 是包含 H 的 G 的最小余格（利用 VI 第 16 页的注记）。
\end{enumerate}

\begin{enumerate}
\def\labelenumi{\arabic{enumi})}
\setcounter{enumi}{12}
\tightlist
\item
  \begin{enumerate}
  \def\labelenumii{\alph{enumii})}
  \tightlist
  \item
    一个幺半群 \(\mathbf{M}\) 称为下半格序的（或简称半格序的），如果它是一个序幺半群，如果 \(\inf(x, y)\) 对所有 \(x, y \in \mathbf{M}\) 存在，并且如果
  \end{enumerate}
\end{enumerate}

\[
\inf (x + z, y + z) = \inf (x, y) + z
\]

对所有 \(x, y, z \in \mathbf{M}\) 。证明恒等式：

\[
\inf (x, z) + \inf (y, z) = \inf (x + y, z + \inf (x, y, z))
\]

\[
\inf (x, y, z) + \inf (x + y, y + z, z + x) = \inf (x, y) + \inf (y, z) + \inf (z, x).
\]

由这些关系中的第一个，推导出 \(x \leqslant z\) 与 \(y \leqslant z\) 蕴含 \(x + y \leqslant z + \inf(x, y)\) 。

证明命题 11 及其推论（VI，第 14-15 页）在半格序幺半群中成立。b) 设 M 为半格序幺半群，N 为子幺半群，使得 \(\inf_{\mathbf{M}}(x,y) = \inf_{\mathbf{N}}(x,y)\) 对 \(x,y\in \mathbf{N}\) 成立，且 N 的元素在 M 中可逆。证明由元素 \(x - y\) （给定 \(x,y\in \mathbf{N}\) ）组成的 M 的子群 G 是格序的。

\begin{enumerate}
\def\labelenumi{\arabic{enumi})}
\setcounter{enumi}{13}
\tightlist
\item
  证明：在半格序幺半群 M 中，
\end{enumerate}

\[
\inf \left(x _ {i} + y _ {i}\right) \geqslant \inf \left(x _ {i}\right) + \inf \left(y _ {i}\right)
\]

对所有包含 n 项的有限序列 \((x_{i})\) , \((y_{i})\) 成立。推导出不等式

\[
(x + y) ^ {+} \leqslant x ^ {+} + y ^ {+}, \quad | x ^ {+} - y ^ {+} | \leqslant | x - y |
\]

在任何格序群中。

\begin{enumerate}
\def\labelenumi{\arabic{enumi})}
\setcounter{enumi}{14}
\tightlist
\item
  证明：
\end{enumerate}

\[
\left| x ^ {+} - y ^ {+} \right| + \left| x ^ {-} - y ^ {-} \right| = | x - y |
\]

在格序群中（注意 \(x - y \leqslant |x - y|\) 与 \(|x| - |y| \leqslant |x - y|\) ）。

\begin{enumerate}
\def\labelenumi{\arabic{enumi})}
\setcounter{enumi}{15}
\tightlist
\item
  证明：
\end{enumerate}

\[
n. \inf (x, y) + \inf (n x, n y) = 2 n. \inf (x, y)
\]

在半格序幺半群中（参见习题 8）。由此推出 \(\inf(nx, ny) = n \cdot \inf(x, y)\) ，如果 \(\inf(x, y)\) 是正则元。

\begin{enumerate}
\def\labelenumi{\arabic{enumi})}
\setcounter{enumi}{16}
\tightlist
\item
  设 \(x, y, z, t\) 是半格序幺半群 \(M\) 的元素，使得 \(z \geqslant 0\) 且 \(t \geqslant 0\) 。证明不等式
\end{enumerate}

\[
\inf (x + z, y + t) + \inf (x, y) \geqslant \inf (x + z, y) + \inf (x, y + t).
\]

\begin{enumerate}
\def\labelenumi{\arabic{enumi})}
\setcounter{enumi}{17}
\item
  设 \((\mathbf{G}_{\iota})_{\iota \in I}\) 是以良序集 \(I\) 为指标的全序群族；在直和 \(\mathbf{G}'\) （诸 \(\mathbf{G}_{\iota}\) 的直和）上定义一个序，取严格正元素为这样的非零 \((x_{\iota})\) ：满足 \(x_{\iota} > 0\) ，其中 \(\iota\) 为使 \(x_{\iota} \neq 0\) 的第一个指标。证明 \(\mathbf{G}'\) 是全序的。
\item
  设 \(G\) 是一个加法群，\((P_{\alpha})\) 是 \(G\) 的非空子集族，使得 \(P_{\alpha} + P_{\alpha} \subset P_{\alpha}\) 且 \(P_{\alpha} \cap (-P_{\alpha}) = \{0\}\) ；设 \(G_{\alpha}\) 表示在 \(G\) 上取 \(P_{\alpha}\) 为正元素之集所得到的序群。证明 \(P + P \subset P\) 且 \(P \cap (-P) = \{0\}\) ，其中 \(P = \bigcap_{\alpha} P_{\alpha}\) ；如果 \(H\) 是在 \(G\) 上取 \(P\) 为
\end{enumerate}

正元素之集所得到的序群，证明 H 同构于诸 \(G_{\alpha}\) 的乘积的对角子群。

\begin{enumerate}
\def\labelenumi{\arabic{enumi})}
\setcounter{enumi}{19}
\tightlist
\item
  设 \(G\) 是一个加法群，且设 \(P\) 是 \(G\) 的满足下列条件的子集：
\end{enumerate}

\begin{enumerate}
\def\labelenumi{\arabic{enumi}.}
\tightlist
\item
  \(\mathrm{P} + \mathrm{P} = \mathrm{P}\) ；2. \(\mathrm{P} \cap (-\mathrm{P}) = \{0\}\) ；3. 对每个整数 \(n > 0\) ，\(nx \in \mathrm{P}\) 蕴含 \(x \in \mathrm{P}\) （条件 (C)）。
\end{enumerate}

\begin{enumerate}
\def\labelenumi{\alph{enumi})}
\item
  如果 \(a\) 是 \(G\) 的元素使得 \(a \notin P\) ，证明存在一个子集 \(P'\) （\(G\) 的一个子集），满足条件 (C)，使得 \(P \subset P'\) 且 \(-a \in P'\) （取 \(P'\) 为由这样的 \(x \in G\) 所成的集合：存在整数 \(m > 0\) 与 \(n \geqslant 0\) 以及元素 \(y \in P\) ，使得 \(mx = -na + y\) 。
\item
  由 a) 推出，P 是 G 的满足下列条件的子集 T 的交集： \(T + T = T\) , \(T \cap (-T) = \{0\}\) , \(T \cup (-T) = G\) （即使 G 成为全序群）且 \(P \subset T\) （利用佐恩引理）。
\item
  特别地，若 G 是一个每个非零元素均无限阶的加法群，则 G 的所有满足 \(T + T = T\) , \(T \cap (-T) = \{0\}\) 且 \(T \cup (-T) = G\) 的子集 T 的交集是 \(\{0\}\) 。
\end{enumerate}

\begin{enumerate}
\def\labelenumi{\arabic{enumi})}
\setcounter{enumi}{20}
\item
  一个序群称为格可序的，如果它同构于某个格序群的子群。证明：一个序群是格可序的，当且仅当对每个整数 \(n > 0\) ，关系 \(nx \geqslant 0\) 蕴含 \(x \geqslant 0\) （为证该条件充分，利用习题 20，b) 与习题 19）。证明每个格可序群都同构于全序群乘积的一个子群。
\item
  设 G 是一个格可序群（视为 Z-模），E 是向量空间 \(G_{(Q)}\) （II，第 277 页）；证明 E 上存在唯一的与 E 的加法群结构相容的序，它诱导 G 上给定的序，并且使 E 是格可序的。
\item
  设 G 是格序群 \(Z \times Z\) （其中 Z 具有通常的序），H 是由 \((2, -3)\) 生成的 G 的凸子群；证明序群 G/H 不是格可序的（参见 VI，第 30 页，习题 4）。
\item
  设 A 是一个交换环，I 是 A 的所有理想组成的集合。则 \(\mathfrak{a}(\mathfrak{b} + \mathfrak{c}) = \mathfrak{a}\mathfrak{b} + \mathfrak{a}\mathfrak{c}\) 对 I 中所有的 a、b、c 成立（I，第 107 页）；换言之，I 连同序关系 \(a \supset b\) 与合成律 \((\mathfrak{a}, \mathfrak{b}) \mapsto \mathfrak{a}\mathfrak{b}\) 是一个半格序幺半群（VI，第 31 页，习题 13），I 中两个元素 a、b 的上确界为 \(a \cap b\) ，它们的下确界为 \(a + b\) 。
\end{enumerate}

\begin{enumerate}
\def\labelenumi{\alph{enumi})}
\tightlist
\item
  设 K 为一个域，且 \(A = K[X, Y]\) 为 K 上两个未定元的多项式环。
\end{enumerate}

\[
\mathbf {a} = (\mathbf {X}), \mathbf {b} = (\mathbf {Y}), \mathbf {c} = (\mathbf {X} + \mathbf {Y})
\]

\[
\begin{array}{c} (\mathfrak {a} \cap \mathfrak {b}) + \mathfrak {c} \neq (\mathfrak {a} + \mathfrak {c}) \cap (\mathfrak {b} + \mathfrak {c}), \\ (\mathfrak {a} + \mathfrak {b}) \cap \mathfrak {b} \neq (\mathfrak {a} \cap \mathfrak {c}) + (\mathfrak {b} \cap \mathfrak {c}), \\ (\mathfrak {a} \cap \mathfrak {b}) (\mathfrak {a} + \mathfrak {b}) \neq (\mathfrak {a} (\mathfrak {a} + \mathfrak {b})) \cap (\mathfrak {b} (\mathfrak {a} + \mathfrak {b})). \end{array}
\]

\begin{enumerate}
\def\labelenumi{\alph{enumi})}
\setcounter{enumi}{1}
\tightlist
\item
  在环 A 中，1 是 X 与 Y 的一个最大公因子，但 \((\mathbf{X}) + (\mathbf{Y}) \neq \mathbf{A}\) 。
\end{enumerate}

\begin{enumerate}
\def\labelenumi{\arabic{enumi})}
\setcounter{enumi}{24}
\tightlist
\item
  设 I 是交换环 A 的理想的半格序幺半群（习题 24）。要使有限同余式组 \(x \equiv a_i(a_i)\) 只要其中任意两个有公共解便有解（也就是说，只要
\end{enumerate}

\[
a _ {i} \equiv a _ {j} (\mathfrak {a} _ {i} + \mathfrak {a} _ {j})
\]

对每对指标 \(i, j)\) 成立），其必要且充分条件是每个合成律

\[
(a, b) \mapsto a \cap b, (a, b) \mapsto a + b,
\]

在 I 中对另一个合成律是分配的（«中国剩余定理»）。可按如下步骤进行：

\begin{enumerate}
\def\labelenumi{\alph{enumi})}
\item
  如果 \((\mathfrak{a}_1 \cap \mathfrak{a}_2) + (\mathfrak{a}_1 \cap \mathfrak{a}_3) = \mathfrak{a}_1 \cap (\mathfrak{a}_2 + \mathfrak{a}_3)\) ，并且三个同余式 \(x \equiv a_i(\mathfrak{a}_i)\) ( \(i = 1, 2, 3\) ) 中任意两个都有公共解，则这三个同余式有公共解（设 \(x_{12}\) 是 \(x \equiv a_1(\mathfrak{a}_1)\) 与 \(x \equiv a_2(\mathfrak{a}_2)\) 的一个公共解，且设 \(x_{13}\) 是 \(x \equiv a_1(\mathfrak{a}_i)\) 与 \(x \equiv a_3(\mathfrak{a}_3)\) 的一个公共解；证明同余式 \(x \equiv x_{12}(\mathfrak{a}_1 \cap \mathfrak{a}_2)\) 与 \(x \equiv x_{13}(\mathfrak{a}_1 \cap \mathfrak{a}_3)\) 有公共解）。
\item
  如果任意由三个同余式组成的组，只要其中任意两个有公共解，便有公共解，证明下列分配律成立
\end{enumerate}

\[
\mathfrak {a} + (\mathfrak {b} \cap \mathfrak {c}) = (\mathfrak {a} + \mathfrak {b}) \cap (\mathfrak {a} + \mathfrak {c}), \quad \text { and } \quad \mathfrak {a} \cap (\mathfrak {b} + \mathfrak {c}) = (\mathfrak {a} \cap \mathfrak {b}) + (\mathfrak {a} \cap \mathfrak {c}).
\]

（对于第一个，注意对所有 \(x \in (\mathfrak{a} + \mathfrak{b}) \cap (\mathfrak{a} + \mathfrak{c})\) 存在 \(y \in \mathfrak{b} \cap \mathfrak{c}\) 使得 \(y \equiv x(\mathfrak{a})\) ；对于第二个，注意对所有 \(x \in \mathfrak{a} \cap (\mathfrak{b} + \mathfrak{c})\) 存在 \(y \in \mathfrak{a} \cap \mathfrak{b}\) 使得 \(y \equiv x(\mathfrak{c})\) 。）注意由 \(a)\) 与 \(b)\) ，第二个分配律蕴含第一个（参见集合论，III，第 217 页，习题 16）。

\begin{enumerate}
\def\labelenumi{\alph{enumi})}
\setcounter{enumi}{2}
\item
  对所考虑的同余式个数作归纳，用与 a) 类似的论证证明«中国剩余定理»。
\item
  * 主理想整环情形的翻译。 *
\end{enumerate}

\begin{enumerate}
\def\labelenumi{\arabic{enumi})}
\setcounter{enumi}{25}
\item
  证明理想 (X) 在多项式环 K{[}X, Y{]} （其中 K 是一个交换域）的理想幺半群中满足 VI 第 17 页命题 14 的条件 (P)，但不是极大的。
\item
  设 A 是以 \((1, e)\) 为基的二次 Z-代数，其中 \(e^{2} = -5\) 。证明 A 是一个整环；在 A 中我们有 \(9 = 3 \cdot 3 = (2 + e)(2 - e)\) ；证明 3、\((2 + e)\) 与 \((2 - e)\) 是 A 的不可约元，但不满足 VI 第 17 页命题 14 (DIV) 的条件。
\item
  设 G 是一个格序群，其正元素之集为 P。P 的两个元素 x 与 y 称为等价的，如果与其中一个互素的任何元素也与另一个互素；此关系下的等价类称为 P 的线程；设 \(\bar{x}\) 表示包含 x 的线程。
\end{enumerate}

\begin{enumerate}
\def\labelenumi{\alph{enumi})}
\item
  设 \(\overline{a}\) 与 \(\overline{b}\) 是线程；设 \(x, x_1\) 是 \(\overline{a}\) 的元素，\(y, y_1\) 是 \(\overline{b}\) 的元素；证明如果与 \(x\) 互素的每个元素也与 \(y\) 互素，则与 \(x_1\) 互素的每个元素也与 \(y_1\) 互素。按此方式在 \(\overline{a}\) 与 \(\overline{b}\) 之间定义一个关系，记作 \(\overline{a} \geqslant \overline{b}\) ；证明这是线程之集 F 的一个序。
\item
  证明 \(\mathbf{F}\) 在此序下是一个格，并且
\end{enumerate}

\[
\inf (\bar {a}, \bar {b}) = \overline {{\inf (a , b)}}, \quad \sup (\bar {a}, \bar {b}) = \overline {{\sup (a , b)}} = \overline {{a + b}}
\]

（利用 VI 第 15 页的推论 3）。证明 F 有一个最小元素，即线程 \(\overline{0} = \{0\}\) 。

\begin{enumerate}
\def\labelenumi{\alph{enumi})}
\setcounter{enumi}{2}
\item
  称两个线程 \(\overline{a}\) , \(\overline{b}\) 互素，如果 \(\inf (\overline{a},\overline{b}) = \overline{0}\) 。证明如果 \(\overline{a}\) 与 \(\overline{b}\) 是两个 \(\neq \overline{0}\) 的线程，且若 \(\overline{a} < \overline{b}\) , 则存在一个线程 \(\overline{c}\) ，与 \(\overline{a}\) 互素，使得 \(\overline{c} < \overline{b}\) 。
\item
  称一个线程 \(\bar{m} \neq \bar{0}\) 不可约，如果它是 \(\neq \bar{0}\) 的线程所成之集的极小元素。证明一个不可约线程是全序的（如果 \(a\) 与 \(b\) 是 \(\bar{m}\) 的两个元素，考虑包含 \(a - \inf(a, b)\) 与 \(b - \inf(a, b)\) 的线程，并利用 \(b\) ）。再证明 \(\bar{m}, -\bar{m}\) 与 \(\{0\}\) 的并集是一个全序子群 \(\mathrm{H}(\bar{m})\) ，即 \(\mathbf{G}\) 的一个子群。
\item
  给定 G 的互不相同的不可约线程的族 \((\bar{m}_{\iota})\) ，证明由这些 \(\bar{m}_{\iota}\) 的并集生成的子群 H 同构于这些群 \(\mathrm{H}(\bar{m}_{\iota})\) 的直和（化归到有限族的情形，并用归纳法论证）。
\item
  证明：在非零全序群的乘积（相应地，直和）中，不可约线程的元素恰是那些除一个坐标（为正的）外其余坐标全为零的元素。由此推出，一个序群至多以一种方式表示为非零全序群的乘积（相应地，直和）（换言之，因子是唯一确定的）。
\end{enumerate}

\begin{enumerate}
\def\labelenumi{\arabic{enumi})}
\setcounter{enumi}{28}
\tightlist
\item
  序群 G 称为完全格序的，如果它是滤过的，且 G 的每个有上界的非空子集在 G 中都有上确界。下半格序幺半群称为完全半格序的，如果 \(\inf(x_{\alpha})\) 对每个非空
\end{enumerate}

的元素族 \((x_{\alpha})\) （属于 M 且有下界的）存在，并且如果

\[
\inf _ {\alpha , \beta} (x _ {\alpha} + y _ {\beta}) = \inf _ {\alpha} (x _ {\alpha}) + \inf _ {\beta} (y _ {\beta})
\]

对任意两个有下界的非空族 \((x_{\alpha})\) , \((y_{\beta})\) 都成立。完全格序群是完全半格序幺半群。

\begin{enumerate}
\def\labelenumi{\alph{enumi})}
\item
  证明：滤过群 G 是完全格序的必要且充分条件是，G 的正元素之集 P 的每个非空子集在 P 中都有下确界。
\item
  完全格序群的任意乘积仍是完全格序的。
\item
  完全格序群的每个滤过凸子群都是完全格序的。
\end{enumerate}

\begin{enumerate}
\def\labelenumi{\arabic{enumi})}
\setcounter{enumi}{29}
\tightlist
\item
  设 G 是一个序群。对 G 的每个子集 A，令 \(m(A)\) （相应地，M(A)）表示 A 在 G 中的下界（相应地，上界）所成的集合；则 \(m(A) = -M(-A)\) 。a) 关系 \(A \subset B\) 蕴含 \(m(B) \subset m(A)\) 与 \(\mathbf{M}(m(A)) \subset \mathbf{M}(m(B))\) 。
\end{enumerate}

\begin{enumerate}
\def\labelenumi{\alph{enumi})}
\setcounter{enumi}{1}
\item
  我们有 \(\mathbf{A} \subset \mathbf{M}(m(\mathbf{A}))\) 与 \(\mathbf{M}(m(\mathbf{M}(\mathbf{A}))) = \mathbf{M}(\mathbf{A})\) 。
\item
  如果 \((\mathbf{A}_{\iota})\) 是 \(G\) 的任意子集族，则 \(m\left(\bigcup_{\iota} \mathbf{A}_{\iota}\right) = \bigcap_{\iota} m(\mathbf{A}_{\iota})\) 。
\item
  每个集合 M(B)，其中 B 是 G 中有上界的非空子集，称为 G 中的主集。对 G 中任意有下界的非空子集 A，集合 M(m(A)) 是包含 A 的最小主集；将其记为 \(\langle A\rangle\) ；若 \(A\subset B\) ，则 \(\langle A\rangle\subset\langle B\rangle\) ；若 A 在 G 中有下确界 a，则 \(\langle A\rangle=M(a)\) ，我们也把它写作 \(\langle a\rangle\) 。
\item
  若 A 与 B 是 G 中两个有下界的非空子集，则 \(\langle A + B \rangle = \langle \langle A \rangle + B \rangle\) 。由此推出，在 G 的主集所成的集合 \(\mathfrak{M}(G)\) 中，映射 \((A, B) \mapsto \langle A + B \rangle\) 是一个结合且交换的合成律，对此律 \(\langle 0 \rangle = P\) 是单位元，序关系 \(A \supset B\) 与该律相容，且若 G 是滤过的，则 \(\mathfrak{M}(G)\) 在此结构下是完全半格序幺半群（习题 29）。此外，映射 \(x \mapsto \langle x \rangle\) 从 G 到 \(\mathfrak{M}(G)\) 是从序群 G 到幺半群 \(\mathfrak{M}(G)\) 的一个子群上的同构。
\item
  如果 \(\mathfrak{M}(G)\) 的一个元素 A 在此幺半群的合成律下可逆，则其逆元为 \(\mathbf{M}(-\mathbf{A}) = -m(\mathbf{A})\) （注意，若 B 是 A 的逆元，则关系 \(\mathbf{B} \subset \mathbf{M}(-\mathbf{A})\) 与 \(\mathbf{A} + \mathbf{M}(-\mathbf{A}) \subset \langle 0 \rangle\) 成立，由此得 \(\langle \mathbf{A} + \mathbf{B} \rangle = \langle \mathbf{A} + \mathbf{M}(-\mathbf{A}) \rangle\) ）。
\item
  要使 \(\mathfrak{M}(G)\) 的一个元素 A 可逆，必要且充分条件是 \(x + A \subset A\) 蕴含 \(x \geqslant 0\) （写出 \(0 \in \langle A + M(-A) \rangle\) 的一个表达式）。
\end{enumerate}

\begin{enumerate}
\def\labelenumi{\arabic{enumi})}
\setcounter{enumi}{30}
\tightlist
\item
  序群 G 称为阿基米德的，如果使得 nx（n 为正整数）所成之集有下界的元素 \(x \in G\) 只有 G 的正元素。
\end{enumerate}

\begin{enumerate}
\def\labelenumi{\alph{enumi})}
\item
  序群 G 中主集所成的幺半群 \(\mathfrak{M}(G)\) 是群，当且仅当 G 是阿基米德的（利用习题 30 的 g) 与 d)）。如果此外 G 还是滤过的，则 \(\mathfrak{M}(G)\) 是完全格序群。
\item
  由此推出，一个滤过群 G 同构于某个完全格序群的子群，当且仅当 G 是阿基米德的。
\end{enumerate}

\begin{enumerate}
\def\labelenumi{\arabic{enumi})}
\setcounter{enumi}{31}
\tightlist
\item
  \begin{enumerate}
  \def\labelenumii{\alph{enumii})}
  \tightlist
  \item
    设 G 是一个完全格序群，H 是 G 的一个子群。对群 H 的每个主子集 A，令 \(x_{\mathbf{A}}\) 表示 A 在 G 中的下确界；证明映射 \(\mathbf{A} \mapsto x_{\mathbf{A}}\) 从 \(\mathfrak{M}(\mathbf{H})\) 到 G 是单射（证明 A 恰是那些 \(y \in \mathbf{H}\) （满足 \(y \geqslant x_{\mathbf{A}}\) ）所成的集合）。
  \end{enumerate}
\end{enumerate}

\begin{enumerate}
\def\labelenumi{\alph{enumi})}
\setcounter{enumi}{1}
\item
  对 H 的每个在 H 中有下界的子集 B，令 \(\langle B\rangle\) 表示由 B 在 H 中生成的主集（\(\mathfrak{M}(H)\) 的一个元素）。如果对 H 的每个在 H 中有下界的子集 B，等式 \(\inf B = \inf \langle B \rangle\) 成立（下确界在 G 中取），证明映射 \(A \mapsto x_{A}\) 是从序群 \(\mathfrak{M}(H)\) 到 G 的一个子群上的同构（参见习题 30，e)）。
\item
  如果 G 的每个元素都是 H 的某个子集的下确界，证明对 H 的每个在 H 中有下界的子集 B，有 inf B = inf \(\langle B\rangle\) （下确界在 G 中取）。（注意，当 inf B ∈ H 时这是成立的；一般情形，考虑一个元素 \(x \in G\) ，使得 \(x + inf B \in H\) ，并利用 x 是 G 的某个子集的下确界这一事实，以及习题 30，e)。
\end{enumerate}

\(d) ^{*}\) 设 \(G\) 是完全格序群 \(\mathbf{Z} \times \mathbf{Z} \times \mathbf{R}\) 。设 \(\theta > 0\) 是一个无理数；设 \(H\) 是由 \((x, y, z)\) （满足 \(\theta(z - x) + y = 0\) ）组成的子群。证明不存在序群 \(\mathfrak{M}(H)\) 到 \(G\) 的任何子群上的、其在 \(H\) 上的限制为恒等映射的同构（注意 \(\mathfrak{M}(H)\) 同构于 \(K = Z \times R\) ；考虑 \(K\) 中由这样的元素 \(u\) 组成的子群：对每个整数 \(n > 0\) 存在 \(v \in K\) 使得 \(nv = u\) ，以及 \(G\) 的类似子群）。 *

\begin{enumerate}
\def\labelenumi{\arabic{enumi})}
\setcounter{enumi}{32}
\tightlist
\item
  \begin{enumerate}
  \def\labelenumii{\alph{enumii})}
  \tightlist
  \item
    一个全序群 G 是阿基米德的，当且仅当对 G 的任意一对元素 \(x > 0\) , \(y > 0\) ，存在整数 \(n > 0\) 使得 \(y \leqslant nx\) 。
  \end{enumerate}
\end{enumerate}

* b) 任何全序的、阿基米德的完全格序群 G 都同构于 \(\{0\}\) 、Z 或 R（排除前两种情形，在 G 中固定一个 a \textgreater{} 0，并使每个 \(x \in G\) 对应于使得 \(qx \leq pa\) 的有理数 p/q 的下确界）。

\begin{enumerate}
\def\labelenumi{\alph{enumi})}
\setcounter{enumi}{2}
\item
  由此推出，每个阿基米德全序群 G 都同构于 R 的一个子群（注意 \(\mathfrak{M}(G)\) 是全序的）。*
\item
  字典序积 \(Z \times Z\) 是全序的且不是阿基米德的。
\end{enumerate}

\begin{enumerate}
\def\labelenumi{\arabic{enumi})}
\setcounter{enumi}{33}
\item
  设 \(G = Z^{N}\) 是可数无限个全序群 Z 的乘积，并令 \(H = Z^{(N)}\) 为滤过凸子群，即 G 的各因子的直和。证明格序群 G/H（VI，第 30 页，习题 4）不是阿基米德的，尽管 G 与 H 都是完全格序的。
\item
  序群 G 中的主集 A 称为有限生成的，如果存在有限集 F 使得 \(A = \langle F \rangle\) 。我们称 G 是半阿基米德的，如果每个有限生成的主集在幺半群 \(\mathfrak{M}(G)\) 中都可逆。每个格序（相应地，阿基米德）群都是半阿基米德的（习题 31，a)）。
\end{enumerate}

\begin{enumerate}
\def\labelenumi{\alph{enumi})}
\item
  证明每个滤过的半阿基米德群都是格序的（若 \(nx \geqslant 0\) ，考虑主集 \(\langle F \rangle\) ，其中 \(F = \{0, x, \ldots, (n - 1)x\}\) ，并利用习题 30 的 g) 以及它可逆这一事实）。
\item
  * 设 K 是字典序积 \(R \times R\) ，而 G 是（通常的）乘积群 \(K \times R\) 。设 \(\theta\) 是满足 \(0 < \theta < 1\) 的无理数，并设 H 是由 \(((1, 0), 0)\) 、\(((\theta, 0), \theta)\) 与 \(((0, x), 0)\) （其中 x 遍历 R）生成的子群。证明 H 作为两个全序群的乘积的子群，不是半阿基米德的（考虑由 H 的两个元素 \(((1,0),0)\) 与 \(((\theta ,0),\theta)\) 生成的主集，并证明它不可逆）。*
\item
  设 G 是一个半阿基米德滤过群；证明 \(\mathfrak{M}(G)\) 中由诸有限生成主集及其逆元生成的子幺半群是一个格序群（参见 VI，第 31 页，习题 13，b)）。
\item
  * 设 \(\theta > 0\) 是一个无理数，并设 \(G\) 是群 \(\mathbf{Z} \times \mathbf{Z}\) ，其中正元素之集 \(P\) 取为 \((x, y)\) （满足 \(x \geqslant 0\) 且 \(y \geqslant \theta x\) ）所成之集。证明 \(G\) 是阿基米德群，但在 \(\mathfrak{M}(G)\) 中，一个非 \(\langle a \rangle\) 形式的有限生成主集的逆元不是有限生成的。*
\end{enumerate}

\begin{enumerate}
\def\labelenumi{\arabic{enumi})}
\setcounter{enumi}{35}
\tightlist
\item
  设 \(\mathbf{G}\) 是一个以 \(\mathbf{P}\) 为正元素之集的滤过群。则 \(\mathbf{G}\) 同构于形如 \(\mathbf{Z}^{(I)}\) 的群，当且仅当存在映射 \(x \mapsto d(x)\) （从 \(\mathbf{P}\) 到 \(\mathbf{N}\) ），使得对所有 \(b\) ，或者 \(b \geqslant a\) ，或者存在一个元素 \(c\) ，它属于主集 \(\langle a, b \rangle\) ，使得 \(d(c) < d(a)\) 。（证明该条件在 \(\mathbf{Z}^{(I)}\) 中成立）由映射 \(d(x_{i}) = \sum_{i} x_{i}\) （当 \(x_{i} \geqslant 0\) 对所有 \(i\) 成立时）满足。
\end{enumerate}

反之，若该条件满足，证明每个主集 \(\mathbf{A} \subset \mathbf{P}\) 具有形式 \(\langle a \rangle\) ，办法是取一个元素 \(a\) 属于 \(\mathbf{A}\) ，使得 \(d(a) \leqslant d(x)\) 对所有 \(x \in \mathbf{A}\) 成立。应用 VI 第 18 页定理 2。）
